% ============================================================
% §4 模型假设与合理性论证 (Stage 04 产物)
% 竞赛: 电工杯 B题 (2026)
% ============================================================

\section{模型假设与合理性论证}

建模本质上是对现实世界的受控简化——假设的质量直接决定模型的解释力边界。本文在充分尊重题目背景与附件数据结构的前提下，提炼以下四条核心假设，每条均附类型分类与支撑依据，避免"假设无支撑"与"假设过多导致模型脆弱"两类常见反模式。

\begin{table}[H]
\centering
\caption{核心模型假设清单}
\label{tab:assumptions}
\begin{tabular}{p{0.8cm}p{5.5cm}p{2.2cm}p{4.5cm}}
\toprule
\textbf{编号} & \textbf{假设内容} & \textbf{类型} & \textbf{支撑依据} \\
\midrule
A1 & 各小区老年人口的马尔可夫状态转移服从独立同质的年际递推规律，即自理$\to$半失能$\to$失能的退化过程在所有10个小区中遵循相同的转移概率矩阵$\mathbf{A}$，且各小区之间不存在老年人口的跨区迁移 & 模型隐含 & 题目明确给出统一转移概率；独立同质假设是马尔可夫链的标准前提，10个小区隶属于同一街道，微观迁移可忽略 \\
\midrule
A2 & 五年预测窗口内，当地工资水平、物价指数与政府养老补贴政策保持平稳，不发生结构性突变 & 环境假设 & 题目明确要求政策稳定；该假设将静态成本参数的有效期从基年合理外推至第5年末，避免引入不可观测的政策冲击 \\
\midrule
A3 & 老人选择满意度的行为完全由题目给定的三因素加权评分函数$S=0.2S_1+0.3S_2+0.5S_3$刻画，不存在未观测的偏好异质性 & 模型隐含 & 该评分函数为题目唯一给定的满意度框架，保证模型的可验证性与可复现性 \\
\midrule
A4 & 服务站的日最大服务能力$Q_k^{\max}$为硬性上限（不可超载），且日均固定管理成本$O_k$不随实际服务人次的波动而变化 & 数据假设 & 题目对三种规模采用确定性上限表述，符合工程设施”额定容量”的行业惯例；固定管理成本对应人员工资、场地租金等刚性支出 \\
\bottomrule
\end{tabular}
\end{table}

\textbf{假设A1的深度论证}：退化方向性——失能不可逆——对应矩阵$\mathbf{A}$的零元位置编码的不可逆性；增量开放性，即年7\%新增注入，使该系统成为”增生型马尔可夫过程”，区别于标准闭合链的递减趋势；空间独立性，即无跨区迁移，使10组递推可并行计算，与”嵌入式”服务的社区内闭环特征一致。

\textbf{假设A3的补充说明}：$S_3$权重0.5反映出题设对”服务可负担性”的重视。Q3定价优化中，在利润率$\le 8\%$约束下通过降低价格提升$S_3$，将评分规则从描述性工具转化为规范性优化目标。
